This work presents \sintonia, a progressive web application (PWA) for acquiring and
interpreting body composition data from an eight-electrode bioelectrical impedance
scale (CF577BLE). Unlike commercial apps, which only display the indices computed by
the device's proprietary algorithms, the system reads the raw impedance directly from
the scale through the Web Bluetooth API and recomputes body composition using open
scientific equations validated against dual-energy X-ray absorptiometry (DXA), such
as those by Wahrlich et al.\ (2026), Janssen et al.\ (2000) and Wang et al.\ (1999).
To mitigate the bias inherent to bioimpedance, the system records the context of each
measurement through a short questionnaire and applies a within-user statistical
correction model, with mixed effects and uncertainty quantification. Natural-language
interpretation is delivered by the virtual assistant Vida, based on the Qwen 2.5 (7B)
Small Language Model running locally via Ollama, which preserves the privacy of health
data and removes the dependence on third-party cloud services. The results indicate
that it is feasible to turn raw bioimpedance data into clear, scientifically grounded
information.
